\chapter {Basic Theories}
\section {DSP Basics}

DSP is about taking information, like from visual images, sound waves or temperature, etc. from the real-world, which are being passed through an Analog-to-Digital converter as binary signals to perform mathematical operations on it to achieve various goals. E.g.: Analyzing and Filter out Noises, Enhancing the signals, Recognition and Generation of Speech, etc. With a Digital-to-Analog converter the result can be. 

Signals: 

Every signal has its own respective parameter that can be displayed on a graph. As it has already been said that signals are information from the real-world, they contain the characteristic of being continuous. It can be imagined as a voltage that varies its output over time. Those signal parameters (e.g., voltage and time) are being quantified during the AD-Conversion to limit the signal of both parameters to discrete values. Depending on the horizontal axis the domain can be differentiated between time domain and frequency domain. As in the case of FAUST, they are very well suited for implementing time-domain algorithms that can be easily represented as block diagrams. 

Quantization: 

Quantization is the process in AD-Conversion of converting a continuous range of values into a finite sequence of discrete values. The reason for that is because it enables the representation of analog signals in a digital format suitable for storage, transmission and computation.  

\section {FAUST Basics}